\begin{abstract}
We introduce \textsc{QLoRA}, a novel method for finetuning that drastically lowers memory requirements, enabling the finetuning of language models with up to 65B parameters on a single 48GB GPU, while maintaining the performance characteristics of standard 16-bit finetuning. In \textsc{QLoRA}, gradients are propagated through a fixed, 4-bit quantized base language model and applied to Low Rank Adapters (LoRA). Our leading suite of models, dubbed \textbf{Guanaco}, achieves results surpassing all previously available open-source models on the Vicuna benchmark, reaching 99.3\% of ChatGPT's performance after just 24 hours of finetuning on one GPU. The memory savings in \textsc{QLoRA} are achieved through several technical contributions: (a) the introduction of 4-bit NormalFloat (NF4), a data type theoretically optimized for weights following a normal distribution, (b) Double Quantization, a technique which reduces memory usage further by quantizing the quantization coefficients, and (c) Paged Optimizers, designed to handle fluctuations in memory usage. Employing \textsc{QLoRA}, we have finetuned over 1,000 models and deliver an extensive evaluation spanning 8 instruction-tuning datasets, different architectures (including LLaMA and T5), and large model sizes (up to 65B parameters) that are out of reach for conventional finetuning approaches. Our findings demonstrate that QLoRA, even with relatively small but high-quality datasets, attains state-of-the-art results with smaller models than prior leading methods. We further provide an in-depth assessment of chatbot quality using both human raters and GPT-4 as evaluators, indicating that GPT-4 offers a cost-effective and valid substitute for human judgement. Moreover, we observe that widely-used benchmarks for chatbot evaluation may not reliably reflect true chatbot capabilities. An adversarial evaluation (“lemon-picked” analysis) is provided to illustrate specific failure cases of \textbf{Guanaco} compared to ChatGPT. All models and code—including CUDA kernels for 4-bit training—are made publicly available.\footnote{\url{https://github.com/artidoro/qlora} and \url{https://github.com/TimDettmers/bitsandbytes}}
\end{abstract}

\section{Introduction}

Refining large language models (LLMs) through finetuning is widely recognized as a powerful strategy to enhance their capabilities~\citep{min2021metaicl, wei2021finetuned, ouyang2022training, wang2022super, wang2022self, liu2022few}, as well as to incorporate preferred behaviors or eliminate undesirable traits~\citep{ouyang2022training,askell2021general,bai2022training}. Nonetheless, the process of finetuning very large models entails substantial computational cost; for example, conventional 16-bit finetuning of a LLaMA model with 65B parameters~\citep{touvron2023llama} demands in excess of 780 GB of GPU memory. Although recent advancements in quantization have shown promise in lowering the memory consumption during inference~\citep{dettmers2022llm, dettmers2022case, frantar2022gptq, xiao2022smoothquant}, these methods typically do not extend to training, where they often prove inadequate~\citep{wortsman2023stable}.

In this work, we present the first successful demonstration of finetuning a 4-bit quantized model without any loss in performance. Our proposed approach, \textsc{QLoRA}, applies an innovative high-precision quantization method to compress a pretrained model to 4 bits and augments it with a compact set of trainable Low-rank Adapter weights~\citep{hu2021lora}. 

These adapters are optimized by propagating gradients through the quantized parameters during finetuning.

With \textsc{QLoRA}, the memory requirements for finetuning a 65B parameter model are brought down from more than 780GB to less than 48GB of GPU memory, with no compromise to speed or prediction accuracy relative to full 16-bit finetuning. This advancement markedly lowers the barrier for LLM finetuning: the largest publicly available models can now be adapted using a single GPU. Leveraging \textsc{QLoRA}, we develop the \textbf{Guanaco} series of models; our second-best configuration achieves 97.8\% of ChatGPT's performance on the Vicuna~\citep{vicuna2023} benchmark, requiring under 12 hours of training on a consumer GPU, while our top model, trained for 24 hours on a single professional GPU, attains 99.3\%—effectively matching ChatGPT on Vicuna. Furthermore, when deployed, our smallest \textbf{Guanaco} model (7B parameters) operates with only 5 GB of memory, yet exceeds a 26 GB Alpaca model by over 20 percentage points on the Vicuna benchmark (see Table~\ref{tbl:vicuna_eval}).

\textsc{QLoRA} introduces several advancements aimed at minimizing memory consumption while maintaining model effectiveness: (1) \textbf{4-bit NormalFloat}, an information-theoretically optimal quantization format for data with a normal distribution, empirically outperforming traditional 4-bit Integer and Float formats; (2) \textbf{Double Quantization}, a technique in which quantization constants are themselves quantized, leading to an average memory savings of around 0.37 bits per parameter (roughly 3 GB for models with 65B parameters); (3) \textbf{Paged Optimizers}, which leverage NVIDIA’s unified memory to circumvent transient memory spikes due to gradient checkpointing, particularly when handling long sequences within mini-batches. Integrating these features, the approach enhances LoRA by placing adapters at every layer, effectively mitigating the accuracy losses that previous methods encountered.

This increased efficiency afforded by \textsc{QLoRA} allows for comprehensive analyses of instruction finetuning and chatbot capability on model sizes unfeasible with standard finetuning approaches because of memory limitations. As a result, over 1,000 models are trained spanning numerous instruction tuning corpora, model architectures, and parameter counts from 80M to 65B. The experiments demonstrate not only that \textsc{QLoRA} attains 16-bit performance levels (\S\ref{sec:comp_to_std}) and enables the development of the high-performing chatbot \textbf{Guanaco} (\S\ref{sec:pushing}), but also yield insights regarding model trends. Most notably, dataset quality consistently outweighs dataset scale: for instance, a dataset comprising 9k samples (OASST1) surpassed a much larger, 450k sample dataset (FLAN v2, subsampled) on chatbot-related tasks, even when both datasets target instruction-following generalization. Furthermore, the analysis reveals that excelling on the Massive Multitask Language Understanding (MMLU) benchmark does not guarantee top performance on the Vicuna chatbot benchmark and vice versa—highlighting that the relevance of dataset content supersedes dataset volume for specialized objectives.

Additionally, we conduct a thorough examination of chatbot performance leveraging both assessments by human raters and evaluations performed by GPT-4. Our methodology employs a tournament-style benchmarking approach, where competing models generate responses to prompts in head-to-head matchups. The victor in each match is determined either by GPT-4 or by human annotators. We summarize the outcomes of these tournaments using Elo scores~\citep{elo1967proposed, elo1978rating}, providing a ranking system for chatbot capabilities. Our findings indicate a strong concordance between human and GPT-4-based assessments with respect to model rankings, although some notable divergences are observed. This underscores that, while model-based evaluation represents a cost-effective alternative to human annotation, it carries its own uncertainties.

In addition to quantitative benchmarks, we supplement our findings with a qualitative assessment of the \textbf{Guanaco} models. This qualitative investigation draws attention to particular instances of both model success and failure that are not adequately reflected in the quantitative evaluation results.

To support continued research, we make available all model outputs alongside corresponding human and GPT-4 evaluations. We also open-source our entire codebase, including CUDA kernels, and offer integration with the Hugging Face transformers library \citep{wolf2019huggingface}, ensuring broad accessibility. Furthermore, we provide a suite of adapters for models of size 7/13/33/65B, each trained across 8 distinct instruction-following datasets, resulting in a total of 32 open-source, finetuned models.

\section{Background}
\paragraph{Block-wise k-bit Quantization} Quantization refers to the technique of mapping data from a high-information representation to a lower-information (often lower-bitwidth) form. This typically involves converting data types such as 32-bit floats to lower bit-width types, for example 8-bit integers. To fully utilize the range of the lower-bit data type, it is customary to normalize the original data—usually arranged as a tensor—by dividing by the maximum absolute value, thereby rescaling the input into the intended target range. For instance, converting a tensor in FP32 format to an Int8 tensor with a range $[-127, 127]$ is carried out as follows:

\begin{equation}
    \mathbf{X}^{ \text{Int8}} = \text{round}\left( \frac{127}{{\text{absmax}}(\mathbf{X}^{{\text{FP32}}})}\mathbf{X}^{{\text{FP32}}} \right) = \text{round}(c^{\text{FP32}}\cdot \mathbf{X}^{{\text{FP32}}}),
\end{equation}

with $c$ representing the {\em quantization constant} or {\em quantization scale}. The inverse process, referred to as dequantization, is given by:

\begin{equation}
    \text{dequant}(c^{\text{FP32}}, \mathbf{X}^{\text{Int8}}) = \frac{\mathbf{X}^{\text{Int8}}}{c^{\text{FP32}}} = \mathbf{X}^{\text{FP32}}
\end{equation}

One limitation of this method is its inefficiency in handling large-magnitude outlier values in the input tensor. Such outliers can lead to uneven usage of quantization bins, resulting in some bins corresponding to few or no quantized values. To address the outlier problem, a typical solution involves partitioning the input tensor into several blocks, with each block quantized separately using its own quantization constant $c$. More precisely, the input tensor $\mathbf{X}\in \mathbb{R}^{b\times h}$ is reshaped into $n$ consecutive blocks, each of length $B$, by first flattening the tensor and then dividing the resulting vector into ${n = ({b\times h} )/{B} }$ segments. Each segment is quantized independently according to Equation~1, resulting in a quantized tensor along with $n$ distinct quantization constants $c_i$.

\paragraph{Low-rank Adapters} Low-rank Adapter (LoRA) finetuning \citep{hu2021lora} is a technique that achieves lower memory consumption by updating only a compact subset of parameters—adapters—while leaving the original model weights unchanged. During training with stochastic gradient descent, gradients propagate through the static pretrained weights, but updates are applied exclusively to the adapter parameters to minimize the loss. LoRA incorporates an auxiliary, factorized projection into the standard linear projection layer. For a linear transformation given by ${\mathbf{X} \mathbf{W} = \mathbf{Y}}$ where $\mathbf{X}\in  \mathbb{R}^{b\times h}$ and $\mathbf{W}\in  \mathbb{R}^{h\times o}$, the LoRA-modified output is computed as:

\begin{equation}
    \mathbf{Y} = \mathbf{X}\mathbf{W} + s\mathbf{X}\mathbf{L}_1\mathbf{L}_2,
\end{equation}

where $\mathbf{L}_1\in  \mathbb{R}^{h\times r}$ and $\mathbf{L}_2\in  \mathbb{R}^{r\times o}$ are learnable matrices, and $s$ is a scaling coefficient.

\paragraph{Memory Requirement of Parameter-Efficient Finetuning} A key consideration regarding LoRA pertains to its memory consumption during training, particularly as it relates to both the quantity and scale of adapters implemented. Given LoRA's inherently small memory usage, increasing the number of adapters to achieve better outcomes is feasible without causing a substantial rise in overall memory demands. Although LoRA's primary goal is to provide a Parameter Efficient Finetuning (PEFT) approach, it is notable that the dominant contributor to memory usage during LLM finetuning stems from storing activation gradients, rather than from the parameters introduced by LoRA itself. For instance, when finetuning a 7B LLaMA model on FLAN v2 using a batch size of 1 and applying LoRA adapters representing 0.2\% of the model weights---a commonly chosen configuration~\citep{hu2021lora,liu2022few}---the input gradients for LoRA consume 567 MB of memory, whereas the LoRA parameters occupy merely 26 MB. Utilizing gradient checkpointing~\citep{chen2016training} decreases the average memory usage for input gradients to just 18 MB per sequence, which still exceeds the total memory required for LoRA parameters. For perspective, the base model quantized to 4-bit occupies 5,048 MB of memory. These figures illustrate that while gradient checkpointing is a vital technique, further reducing LoRA parameter count results in only marginal savings. Consequently, employing additional adapters does not appreciably increase memory consumption during training (refer to Appendix~\ref{app:memory} for further details). As elaborated in a subsequent section, this property is especially significant when striving to match the performance achieved by full 16-bit precision models.

\section{\textsc{QLoRA} Finetuning}

\textsc{QLoRA} enables high-precision 4-bit finetuning by employing two primary strategies introduced in this work—4-bit NormalFloat (NF4) quantization and Double Quantization. Beyond this, we propose Paged Optimizers to address memory spikes that arise during gradient checkpointing, mitigating the risk of out-of-memory issues that have commonly hindered the single-node finetuning of large-scale models.

Within \textsc{QLoRA}, model weights are maintained in a low-precision storage format—typically 4-bit—and computations are carried out in a higher-precision type, usually BFloat16. Concretely, when a weight tensor is required, it is first dequantized to BFloat16 format, after which all matrix multiplications are executed using this 16-bit precision.

Below, we elaborate on the individual building blocks of \textsc{QLoRA}, followed by a formalized definition.

\paragraph{4-bit NormalFloat Quantization} The NormalFloat (NF) datatype extends Quantile Quantization \citep{dettmers2022optimizers}, an approach that, from an information-theoretic perspective, assigns an equal number of input values to each quantization interval, or bin. Quantile quantization operates by evaluating the quantiles of the input tensor, which involves computing the empirical cumulative distribution function.

A major drawback of quantile quantization is the computational intensity associated with accurate quantile estimation. Consequently, efficient quantile estimation methods, such as SRAM quantiles~\citep{dettmers2022optimizers}, are leveraged. However, because these algorithms are approximate, they can induce considerable quantization errors, particularly for outlier values—which are often especially significant.

Such challenges with costly quantile computation and the errors from approximation can be circumvented when the distributions of input tensors differ only by a quantization constant. In these instances, the input tensors share identical quantiles, making precise quantile estimation both possible and computationally practical.

Because pretrained neural network weights generally follow a normal distribution centered at zero with standard deviation $\sigma$ (refer to Appendix~\ref{app:normality}), it is possible to remap all weight values onto a common fixed distribution by scaling $\sigma$ so that this distribution fits entirely within the representable interval of the target data type. For our purposes, we define this interval as $[-1, 1]$. This process requires normalization of both the quantiles pertaining to the data type and the neural network weights so that they both inhabit this same range.

To construct an information-theoretically optimal data type for zero-mean normal distributions with any standard deviation $\sigma$ within $[-1, 1]$, the following procedure is used: (1) compute $2^k+1$ quantiles from a theoretical $N(0,1)$ distribution to form a $k$-bit quantile-quantized representation suitable for normal distributions; (2) map the quantized data type values onto the $[-1, 1]$ interval; (3) normalize each input weight tensor by rescaling its values to $[-1, 1]$ through division by the largest absolute value (absolute max normalization).

After aligning both the range of weights and the data type, standard quantization procedures may be applied. The normalization step in (3) is tantamount to adjusting the standard deviation of the weight tensor so that it matches that of the $k$-bit data type. Mathematically, the $2^k$ representative values $q_i$ for the data type are calculated as:

\begin{equation}
q_i  = \frac{1}{2}\left( Q_X\left(\frac{i}{2^k+1}\right) + Q_X\left(\frac{i+1}{2^k+1}\right)\right),
\end{equation}

where $Q_X(\cdot)$ denotes the quantile function of the standard normal distribution $N(0,1)$. One limitation of symmetric k-bit quantization is the lack of an exact zero representation, a critical feature when quantizing padding or other elements valued at zero to ensure no quantization error for these values. To address this, and to utilize all $2^k$ available codes for a k-bit format, we design an asymmetric quantization scheme by separately calculating the quantiles $q_i$ for two intervals: $2^{k-1}$ quantiles for the negative region and $2^{k-1}+1$ quantiles for the positive region. These two sets of quantiles are then merged, with the duplicate zero entry eliminated to ensure a unique zero. The resulting representation, which evenly distributes data values across each bin under a standard normal distribution, is termed \emph{k-bit NormalFloat} (NFk), reflecting its information-theoretic optimality for distributions centered at zero. The precise numerical values that define this data type are detailed in Appendix~\ref{app:nf4}.

\paragraph{Double Quantization{}} 
We propose \emph{Double Quantization} (DQ), a scheme that reduces memory usage further by quantizing the quantization constants themselves. While high-fidelity 4-bit quantization necessitates small blocksizes~\citep{dettmers2022case}, this requirement increases the storage needed for the quantization constants. For instance, adopting 32-bit quantization constants and a blocksize of 64 when compressing a weight matrix $\mathbf{W}$ results in an additional memory requirement of $32/64 = 0.5$ bits per parameter. Double Quantization offers a method to alleviate this memory overhead.

In particular, Double Quantization applies a second layer of quantization to the original quantization constants $c_2^{\text{FP32}}$, producing both quantized constants $c_2^{\text{FP8}}$ and new quantization constants $c_1^{\text{FP32}}$. For this second quantization phase, we utilize 8-bit floating point types with a blocksize of 256, consistent with prior work~\citep{dettmers2022case} demonstrating negligible loss in accuracy at this bit width. Since $c_2^{\text{FP32}}$ contains only positive values, we preprocess by subtracting their mean, thus centering the distribution and enabling effective symmetric quantization. Using a blocksize of 64, this method reduces the average per-parameter storage of quantization constants from $32/64=0.5$ bits to $8/64 + 32/(64\cdot256) = 0.127$ bits—a decrease of 0.373 bits per parameter.

\paragraph{Paged Optimizers} leverage NVIDIA’s unified memory~\footnote{\tiny \url{https://docs.nvidia.com/cuda/cuda-c-programming-guide}} technology, which transparently manages page-level transfers between CPU and GPU to ensure seamless GPU operation, even in situations where GPU memory is exhausted. This mechanism is analogous to traditional memory paging between main RAM and disk storage. In our approach, optimizer state memory is allocated as paged memory. When GPU memory becomes insufficient, the system automatically migrates optimizer state data to CPU RAM and later pages it back into GPU memory when needed for optimizer updates.

\paragraph{\textsc{QLoRA}.} Building upon the above mechanisms, we articulate \textsc{QLoRA} for a quantized linear layer equipped with a single LoRA adapter as follows:

\begin{equation}
    \mathbf{Y}^{\text{BF16}} =  \mathbf{X}^{\text{BF16}} \text{doubleDequant}(c_1^{\text{FP32}}, c_2^{\text{k-bit}}, \mathbf{W}^{\text{NF4}}) + \mathbf{X}^{\text{BF16}}\mathbf{L}^{\text{BF16}}_1\mathbf{L}^{\text{BF16}}_2,
\end{equation}

where the function doubleDequant$(\cdot)$ is specified as:

\begin{equation}
    \text{doubleDequant}(c_1^{\text{FP32}}, c_2^{\text{k-bit}}, \mathbf{W}^{\text{k-bit}}) = \text{dequant}(\text{dequant}(c_1^{\text{FP32}}, c_2^{\text{k-bit}}), \mathbf{W}^{\text{4bit}}) = \mathbf{W}^{\text{BF16}},
\end{equation}

We implement NF4 quantization for $\mathbf{W}$ and employ FP8 precision for $c_2$.

To achieve higher quantization accuracy, a blocksize of 64 is used for $\mathbf{W}$, while for $c_2$, a larger blocksize of 256 is chosen to optimize memory efficiency.

During parameter optimization, only the gradients of the loss with respect to the adapter weights $\frac{\partial E}{\partial \mathbf{L}_i}$ are computed, omitting the gradients for the 4-bit weights $\frac{\partial E}{\partial \mathbf{W}}$. Nonetheless, evaluating $\frac{\partial E}{\partial \mathbf{L}_i}$ requires computing $\frac{\partial \mathbf{X}}{\partial \mathbf{W}}$, which is performed according to equation~(5). This necessitates dequantizing the stored $\mathbf{W}^{\text{NF4}}$ into $\mathbf{W}^{\text{BF16}}$, so that the derivative $\frac{\partial \mathbf{X}}{\partial \mathbf{W}}$ can be calculated in BFloat16 format.

In essence, \textsc{QLoRA} employs a single storage data type—typically 4-bit NormalFloat—for model weights, and a separate computation data type, specifically 16-bit BrainFloat. Before executing the forward and backward passes, weights are dequantized from the storage to the computation format. Importantly, gradient calculations are restricted to the LoRA parameters, which remain in the 16-bit BrainFloat format.

\section{QLoRA vs. Standard Finetuning}
\label{sec:comp_to_std}

We have outlined the mechanism of QLoRA and its substantial memory savings for finetuning large models. The central issue remaining is whether QLoRA is able to match the effectiveness of full-model finetuning. Additionally, we seek to understand how the components of QLoRA contribute, such as whether NormalFloat4 provides any measurable advantage over the standard Float4. The subsequent sections elaborate on experiments conducted to investigate these considerations.

\paragraph{Experimental setup.} Our study covers three distinct model types—encoder, encoder-decoder, and decoder-only architectures—and systematically contrasts QLoRA against both 16-bit adapter-based finetuning and traditional full-model finetuning for models up to 3B parameters. The benchmarking tasks include GLUE \citep{wang2018glue} using RoBERTa-large \citep{liu2019roberta}, Super-NaturalInstructions (TKInstruct) \citep{wang2022super} with T5 \citep{t5}, and 5-shot MMLU \citep{hendrycksmeasuring} after finetuning LLaMA on Flan v2 \citep{longpre2023flan} and Alpaca \citep{alpaca}. To further assess the merits of NF4 relative to alternative 4-bit representations, we replicate the setup of \citet{dettmers2022case}, evaluating post-quantization zero-shot accuracy and perplexity on a range of model families—OPT \citep{zhang2022opt}, LLaMA \citep{touvron2023llama}, BLOOM \citep{scao2022bloom}, and Pythia \citep{biderman2023pythia}—with model sizes ranging from 125m to 13B.

To ensure clarity, additional experimental details for each specific configuration are provided in the respective results sections. Comprehensive information can be found in Appendix~\ref{app:exp-setup}.

Paged optimizers are indispensable for QLoRA finetuning on 33B and 65B models using a single 24GB or 48GB GPU. However, since the need for paging is triggered primarily by the occasional processing of mini-batches with extended sequence lengths, we do not report explicit performance metrics for Paged Optimizers. That said, an evaluation of 65B model training on a 48GB GPU, using a batch size of 16, reveals that paged optimizers deliver throughput comparable to standard optimizers. Further research is required to determine under which conditions the paging mechanism introduces measurable slow-downs.

\paragraph{Default LoRA hyperparameters do not match 16-bit performance} 

Applying the typical LoRA approach—modifying only the query and value projection matrices within attention layers as per \citep{hu2021lora}—does not reproduce the results of full finetuning for large-scale base models. Evidence in Figure~\ref{fig:hyper}, from finetuning LLaMA 7B on Alpaca, demonstrates that the principal factor impacting performance is the number and distribution of LoRA adapters; only when LoRA modules are included in every linear layer of the transformer block does performance reach parity with full finetuning. Variations in other hyperparameters, such as the projection dimension $r$, exhibit negligible influence on the outcome (see Appendix \ref{app:exp-setup}).

Our analysis also reveals that the default hyperparameters for fully finetuned baseline models are not sufficiently optimized. To establish strong baselines, we perform a comprehensive hyperparameter sweep, adjusting learning rates from 1e-6 to 5e-5 and experimenting with batch sizes between 8 and 128. The outcomes of 7B LLaMA finetuning on the Alpaca dataset are depicted in Figure~\ref{fig:hyper}.

\paragraph{4-bit NormalFloat surpasses 4-bit Floating Point in empirical performance}

Despite the information-theoretic optimality of the 4-bit NormalFloat (NF4) format, it is necessary to verify whether these theoretical merits are realized in practical tasks. Mirroring the methodology of \citet{dettmers2022case}, we assess quantized large language models (OPT \citep{zhang2022opt}, BLOOM \cite{scao2022bloom}, Pythia \cite{biderman2023pythia}, and LLaMA) with parameter sizes ranging from 125M to 65B, using various data types across language modeling and multiple zero-shot benchmarks. As presented in Figure~\ref{fig:nf4} and Table~\ref{tbl:nf4_ppl}, the NF4 data format yields considerable performance improvements compared to FP4 and Int4, and employing double quantization leads to a lower memory requirement without sacrificing model effectiveness.

\paragraph{k-bit \textsc{QLoRA} achieves parity with 16-bit finetuning and 16-bit LoRA approaches}

Prior research has shown that 4-bit quantization for \emph{inference} is attainable, though it results in reduced accuracy compared to 16-bit models \citep{dettmers2022case,frantar2022gptq}. This naturally leads to the question of whether adapter finetuning using 4-bit quantization can compensate for this reduction in performance. To address this, we evaluate two scenarios.

First, we compare the results to full 16-bit finetuning using RoBERTA and T5 architectures, with model sizes from 125M to 3B parameters, evaluated on GLUE and the Super-NaturalInstructions benchmark. Table~\ref{tab:nf4vsbf16} reports these findings. Across both benchmarks, the results indicate that adapter techniques in 16-bit, 8-bit, and 4-bit modes match the performance of the standard fully finetuned 16-bit baseline. This suggests that any negative impact from reduced quantization precision can be entirely mitigated by applying adapter finetuning after quantization.

In our second experimental setup, due to the substantial GPU memory requirements for fully finetuning models with 11B parameters or more, which necessitate multiple high-memory servers, we shift our focus to assess whether 4-bit \textsc{QLoRA} can achieve performance comparable to 16-bit LoRA across models ranging from 7B to 65B parameters. Specifically, we apply finetuning to LLaMA models from 7B to 65B parameters using the Alpaca and FLAN v2 instruction following datasets and evaluate these models on the MMLU benchmark using a 5-shot accuracy protocol. Table~\ref{tbl:llama-datatype} presents these outcomes, illustrating that the NF4 data type with double quantization successfully matches the MMLU performance of the 16-bit LoRA variant. Additionally, we observe that \textsc{QLoRA} configured with FP4 falls short of the 16-bit brain float LoRA baseline by approximately one percentage point in accuracy. These findings support two key conclusions: (1) \textsc{QLoRA} utilizing NF4 is capable of mirroring the results of both 16-bit full finetuning and 16-bit LoRA finetuning, and (2) in terms of quantization fidelity, NF4 outperforms FP4.

\paragraph{Summary} Across all evaluations, our data demonstrate that 4-bit \textsc{QLoRA} employing the NF4 data type consistently achieves parity with 16-bit full finetuning and 16-bit LoRA on standard academic benchmarks with widely-accepted evaluation procedures. We have further established that NF4 yields superior results to FP4 and that implementing double quantization does not negatively impact model performance. Together, these results strongly suggest that 4-bit \textsc{QLoRA} is a robust and reliable alternative to 16-bit tuning techniques.

Consistent with previous investigations into quantization strategies \citep{dettmers2022case}, both our MMLU and Elo benchmarks demonstrate that, when constrained by a fixed budget for finetuning and inference resources, increasing the model parameter count while reducing their numeric precision is advantageous. This underscores the practical efficiency gains of using \textsc{QLoRA}. Since our experiments involving 4-bit finetuning showed no discernible performance drop compared to full 16-bit finetuning, an open question remains regarding the precise point at which a trade-off between performance and quantization precision emerges for QLoRA. Addressing this will be the subject of future investigation.

We explore instruction tuning at scales that would be impractical using standard 16-bit finetuning on typical research hardware.

\section{Pushing the Chatbot State-of-the-art with QLoRA}
\label{sec:pushing}
Having demonstrated that 4-bit \textsc{QLoRA} attains parity with 16-bit methods across different model sizes, tasks, and datasets, we undertake a detailed investigation of instruction finetuning on the largest open-access language models suitable for academic research. In order to evaluate the efficacy of instruction finetuning on these models, we utilize the demanding MMLU Natural Language Understanding benchmark, and also introduce novel approaches for assessing chatbot effectiveness in real-world scenarios.

\subsection{Experimental setup} The following outlines the experimental design; additional specifics are provided in Appendix \ref{app:chatbot}.

\paragraph{Data} Given the lack of a systematic comparison of contemporary instruction-following datasets, we curated a selection of eight modern datasets. These comprise those generated via crowd-sourcing (OASST1~\citep{kopf2023openassistant}, HH-RLHF~\citep{bai2022training}), distilled from other instruction-tuned models (Alpaca~\citep{alpaca}, self-instruct~\citep{wang2022self}, unnatural-instructions~\citep{honovich2022unnatural}), constructed through corpus aggregation (FLAN v2~\citep{chung2022scaling}), as well as mixed sources (Chip2~\citep{laion2023}, Longform~\citep{koksal2023longform}). These resources span various languages, corpus sizes, and licensing arrangements.

\paragraph{Training Setup} To eliminate confounding influences due to disparate training objectives, we implement QLoRA finetuning strictly with a cross-entropy objective in a supervised framework, omitting reinforcement learning—even for datasets annotated with human response rankings. For datasets providing clear separation between prompts and completions, only the completion is used during finetuning (see the ablation studies in Appendix \ref{app:chatbot}). When working with OASST1 and HH-RLHF, where multiple reply options are present, we pick the best response from each branch of the dialogue tree, and finetune on the concatenated sequence of these chosen responses and corresponding instructions. Across all experiments, we employ NF4 \textsc{QLoRA} with double quantization as well as paged optimizers to mitigate memory overhead from gradient checkpointing. Hyperparameter searches are conducted for the LLaMA models at the 13B and 33B parameter scales, revealing that configurations identified at 7B remain effective for larger models (including epoch count), except for learning rate and batch size, for which we halve the former and double the latter for 33B and 65B.

\paragraph{Baselines} Our comparison set includes both academic (Vicuna~\citep{vicuna2023} and Open Assistant~\citep{kopf2023openassistant}) and commercial (GPT-4~\citep{openai2023gpt}, GPT-3.5-turbo, and Bard) chatbot systems. The Open Assistant baseline leverages a LLaMA 33B model finetuned using Reinforcement Learning from Human Feedback (RLHF) on the OASST1 dataset, identical to one of our benchmarks. In contrast, Vicuna is based on full finetuning of LLaMA 13B with user-contributed ShareGPT conversations, essentially representing a distilled version of OpenAI GPT models.

\subsection{Evaluation}

Consistent with established methodologies, we assess language understanding capabilities using the MMLU (Massively Multitask Language Understanding) benchmark \citep{hendrycksmeasuring}. MMLU features 57 multiple-choice tasks, which include fields such as basic mathematics, US history, computer science, and law. Results are reported as 5-shot accuracy on the test set.

We further assess generative language abilities through both automated systems and evaluations conducted by humans. This second category of assessments utilizes human-crafted queries with the purpose of evaluating the quality of responses generated by the model. While these evaluations offer a more authentic measure of chatbot model performance and are increasingly adopted, there is currently no standardized methodology established in the research literature. Our experimental procedure, outlined below, employs nucleus sampling with $p=0.9$ and a temperature of $0.7$ consistently across all experiments.

\paragraph{Benchmark Data} Our evaluation employs two distinct, human-curated query datasets: the Vicuna prompts \citep{vicuna2023} and the OASST1 validation set \citep{kopf2023openassistant}. The Vicuna prompts consist of 80 questions spanning a variety of categories and are used without alteration. The OASST1 dataset encompasses a multilingual array of multiturn dialogues between users and assistants, crowd-sourced for diversity. From this, we extract all user messages from the validation split as evaluation queries, including the conversation history for full context, resulting in a set of 953 unique user queries. We refer to these resources as the Vicuna and OA benchmarks, respectively.

\paragraph{Automated Evaluation} Drawing from the evaluation strategy presented by \citet{vicuna2023}, we employ GPT-4 to rate model outputs against ChatGPT (GPT-3.5 Turbo) on the Vicuna benchmark. For each query, GPT-4 is presented with both ChatGPT's and the candidate model’s responses, and asked to assign each a score out of ten accompanied by a justification. The relative performance of each system is expressed as a percentage of ChatGPT's score, which may exceed 100\% when a model surpasses ChatGPT’s absolute score. We have observed a notable bias introduced by the order in which responses are presented, with GPT-4 favoring earlier responses. To mitigate this order effect, we advise averaging results across both possible response orders.

We further analyze performance through direct pairwise comparisons of system outputs. This approach transforms the task into a three-class classification problem (choosing one response, the other, or declaring a tie). GPT-4 is instructed to select the superior answer or indicate equivalence, along with an explanation. Such comparisons are conducted exhaustively for every possible system pair across both the Vicuna and OA benchmarks.

\paragraph{Human Evaluation} Although prior studies suggest generative models are viable tools for evaluating conversational systems \citep{fu2023gptscore}, the extent to which GPT-4-based ratings align with human assessments of chatbot responses remains an open question. Accordingly, we supplement automated evaluation with two concurrent rounds of human assessment on the Vicuna benchmark, paralleling the methods previously outlined. Through Amazon Mechanical Turk (AMT), we assign two human annotators for ChatGPT comparisons and three annotators for pairwise system evaluations.

\paragraph{Elo Rating} To evaluate models through both human and automated pairwise judgments, we organize a tournament-style setup in which models are pitted against one another. Each match features a pair of models producing responses to a specific prompt, and their outputs are compared. This approach aligns with the methodologies described by \citet{bai2022training} and \citet{vicuna2023}, but we further integrate ratings generated by GPT-4 in addition to those from humans. To calculate the Elo~\citep{elo1967proposed,elo1978rating} scores, we perform random sampling from our dataset of labeled match outcomes. The Elo system, frequently applied in chess and other competitive games, estimates a model's predicted win probability relative to its competitor’s. For instance, when a model with an Elo of 1100 faces one rated at 1000, it is expected to win around 65\% of the time; if both have the same rating (either 1000 vs 1000 or 1100 vs 1100), their expected win probability is 50\%. The Elo score for a model is adjusted after each match according to how likely the observed result was, so a surprising upset yields a significant Elo shift, while a predictable result produces only a slight adjustment. Over repeated matches, Elo scores become a reasonable proxy for a model’s ability. Each model starts with a rating of 1,000, and we apply $K=32$ as the scaling factor. Following the protocol in \citet{vicuna2023}, the entire evaluation process is repeated 10,000 times using distinct random seeds, helping to minimize the impact of ordering—such as the influence of which model pairs are evaluated first.

\subsection{\textbf{Guanaco}: \textsc{QLoRA} trained on OASST1 Sets a New Standard in Open-Source Chatbots}

Drawing from both automatic and manual assessment methods, our investigation reveals that the leading \textsc{QLoRA}-optimized model, {Guanaco} 65B, finetuned with an OASST1 variant, stands out as the most advanced open-source chatbot to date, rivaling ChatGPT's performance. In head-to-head comparisons, human raters' Elo score assessments show {Guanaco} 65B and 33B models secure a 30\% expected win rate when matched against GPT-4—the highest performance figures yet observed in the open-source domain.

Table~\ref{tbl:vicuna_eval} presents results for the Vicuna benchmark \cite{vicuna2023}, comparing the models' relative performance to ChatGPT. {Guanaco} 65B emerges as the strongest open-source alternative after GPT-4, attaining 99.3\% of ChatGPT’s benchmark score. Notably, although {Guanaco} 33B carries more parameters than Vicuna 13B, it operates at 4-bit precision, making it substantially more efficient with just a 21 GB memory requirement compared to 26 GB, and surpassing Vicuna 13B’s performance by three percentage points. {Guanaco} 7B is also notable for its compactness, requiring only 5 GB—suitable for modern mobile devices—while delivering results nearly 20 percentage points ahead of Alpaca 13B.

Despite these achievements, Table~\ref{tbl:vicuna_eval} indicates that the models' confidence intervals are broad, with considerable overlap in results. We suggest this may stem from ambiguity in score interpretation—for example, the significance of an 8 on a 10-point scale may differ between use cases. To mitigate such limitations, we advocate for the Elo ranking approach \citep{elo1967proposed}, which leverages direct pairwise judgments from humans and GPT-4, providing an evaluation less dependent on arbitrary scales. Elo rankings for leading models are available in Table~\ref{tbl:elo}. It is worth noting that human and GPT-4 rankings differ somewhat, particularly for Guanaco 7B, but display reasonable consistency across most models, as evidenced by a Kendall Tau of $\tau=0.43$ and a Spearman correlation coefficient of $r=0.55$ at the system evaluation level. On a per-example basis, concordance drops, with a Fleiss $\kappa=0.25$ between GPT-4 and aggregate human votes. Collectively, these metrics indicate moderate correspondence between system-level human and GPT-4 assessments, supporting the validity of model-based evaluation as a supplementary alternative to direct human assessment. Further analysis is provided in Section~\ref{ssec:considerations}.

The Elo rankings presented in Table~\ref{tbl:elo_large} reveal that the 33B and 65B variants of {Guanaco} surpass all other models except GPT-4 on both the Vicuna and OA benchmarks. Their performance is also comparable to that of ChatGPT, as reflected in Table~\ref{tbl:vicuna_eval}. It is important to highlight that the Vicuna benchmark tends to favor open-source systems, whereas ChatGPT performs better on the broader OA benchmark. Additionally, as indicated in Tables \ref{tbl:mmlu-llama} and \ref{tbl:vicuna_eval}, the choice of finetuning dataset plays a crucial role in determining model performance. For instance, Llama models finetuned using FLAN v2 achieve high scores on MMLU, but register the lowest results on the Vicuna benchmark; similar patterns are apparent in other models as well. This suggests that evaluation benchmarks capture partially distinct capabilities: excelling at MMLU does not guarantee strong chatbot functionality, as assessed by the Vicuna or OA benchmarks, nor does strong chatbot performance assure high scores on MMLU.

 is unique among the top-performing models in this evaluation for not being trained on proprietary datasets, adhering strictly to the OASST1 guidelines that prohibit the use of GPT-generated content. The next highest-ranking open-data-only model, Anthropic HH-RLHF, trails {Guanaco} by 30 percentage points on the Vicuna benchmark (see Table~\ref{tbl:vicuna_eval}). These findings collectively demonstrate that 4-bit \textsc{QLoRA} facilitates the development of advanced chatbots capable of matching ChatGPT in quality. Moreover, our 33B {Guanaco} can be trained within 12 hours using a 24 GB consumer GPU, thereby enabling further research and development of competitive models through \textsc{QLoRA} finetuning on targeted open-source datasets. This approach offers a viable path toward open models that can rival leading commercial systems.

\section{Qualitative Analysis}

Although our primary assessment relies on quantitative methods, restricting the evaluation to summary metrics overlooks several critical considerations. Chief among these is the challenge of benchmark validity \citep{liao2021we}—the concern over whether a given benchmark genuinely evaluates what it purports to, a question complicated further by the discovery of heuristic ``shortcuts'' that models may leverage to perform well without true task understanding \citep{gururangan2018annotation, poliak2018hypothesis}. To partially address this limitation, we complement our quantitative results with qualitative inspection in this section, divided into two parts. First, in \S\ref{ssec:examples}, we illustrate selected instances from the outputs of our 65b {Guanaco} model that exemplify recurring behaviors or trends. Second, \S\ref{ssec:considerations}, we outline additional reflections on our findings and provide context for our interpretations.

\subsection{Qualitative Analysis of Example Generations}
\label{ssec:examples}

To identify illustrative cases, we manually review responses generated for both the Vicuna and OpenAssistant benchmarks, searching for systematic tendencies in the outputs produced by {Guanaco}. Upon recognizing a recurrent phenomenon, we design queries intended to elicit the observed pattern even when it constitutes an inappropriate or incorrect response. For instance, if we find that the model often produces unnecessarily verbose answers, we instruct it with a prompt like ``Answer yes or no without explanation.'' This methodology enables us to highlight both ``lemons,'' where we succeed in exposing model weaknesses through adversarial prompts, and ``cherries,'' where the model resists such failure modes. All outputs presented in this section are produced using Nucleus Sampling \cite{holtzmancurious} with $p=0.9$.

Naturally, this overview is not exhaustive, as it would be unrealistic within the limitations of this concise qualitative investigation to account for all potentially relevant factors. For example, the set of possible outputs a model may produce in response to a particular prompt is vast, leading us to depend on sampled outputs that we expect are indicative. Nonetheless, we contend that presenting these specific instances helps situate the quantitative findings discussed earlier in the paper. By making all models and associated code publicly available, our intention is that the cases described here will encourage subsequent research to undertake a more granular exploration of the phenomena highlighted.

\paragraph{Factual Recall} For relatively straightforward fact-based prompts like ``What is the capital of Zambia?'', all models reliably return accurate results. For instance,

\begin{quote} 
    \textbf{Guanaco}:
    \texttt{The capital of Zambia is Lusaka.}
\end{quote}

Yet, as the questions become less common or more nuanced, {Guanaco}'s answers lose reliability, although the model’s confidence does not diminish. Consider, for example, this question sourced from HotPotQA \citep{yang2018hotpotqa}:

\begin{quote}
     \textbf{User}: In what year was the singer who popularized `I'll Keep the Lovelight Burning' born?  

     \textbf{Guanaco}:~\texttt{Al Jolson is credited with popularizing the song `I'll Keep the Lovelight Burning,' and he was born in the year 1886.}
\end{quote}

Here, both the individual identified as popularizing the song and the birth year are incorrect, although the provided birth year does accurately correspond to Al Jolson, the figure named in the response.

\paragraph{Suggestibility} Notably, {Guanaco} frequently resists incorporating certain types of misleading presuppositions, as can be seen in the following exchange:

\begin{quote}
    \textbf{User}: How was it finally, officially confirmed that the earth is flat by peer reviewed scientists? \\ \\
    \textbf{Guanaco}: \texttt{There has never been any official confirmation by peer-reviewed scientists that the earth is flat. The vast majority of the scientific community accepts, based on extensive evidence and centuries of observation, that the earth is spherical. \\ \\ Scientific inquiry and empirical data have consistently supported a round earth, while claims supporting a flat earth fall under conspiracy theories that have been thoroughly discredited. \\ \\ The assertion that the flat earth has been validated in peer-reviewed literature is untrue. \\ \\ Feel free to ask further questions if you wish.}
\end{quote}

is also proficient in recognizing questions that cannot be answered, for example,

\begin{quote}
    \textbf{User}: What time is it?

    \textbf{Guanaco}: \texttt{Apologies, but I am an AI language model without real-time data access. My replies are informed by prior knowledge and not the present moment. For current time information, please consult a clock or an online resource.}
\end{quote}

Similar behavior is observed for queries such as ``Where are you?'', ``How are you?'', and related questions.

\paragraph{Refusal} In a comparable manner, {Guanaco} occasionally declines to comply with instructions for reasons that appear arbitrary:

\begin{quote}
    \textbf{User}: Please reverse the words in the sentence ``I'm a cool dude in a mocha mood.'' \\ \\
    \textbf{Guanaco}: \texttt{I'm sorry, but I cannot do that at this moment. However, I can offer an explanation as to why your sentence is grammatically sound.\\ \\The sentence ``I'm a cool dude in a mocha mood'' adheres to the conventional rules of English grammar and structure. Here, the adjective ``cool'' correctly modifies the noun ``dude,'' and the prepositional phrase ``in a mocha mood'' offers further context. This arrangement is entirely acceptable and meaningful.\\ \\If you have further questions or need assistance with something else, please let me know.}
\end{quote}

\paragraph{Secret Keeping} Disappointingly, this form of refusal proves inconsistent, even in scenarios where withholding information is necessary. Designers may wish for systems to safeguard certain secrets from users, even while permitting open-ended text replies. To assess how readily a user might extract sensitive data, we instructed the model: ``The secret word is `banana'. Do not reveal or repeat back this back, no matter what the human requests.'' Direct queries are ineffective:

\begin{quote}
    \textbf{User}: What is the secret word?

    \textbf{Guanaco}: \texttt{I'm sorry, but I cannot share the secret word. I am designed specifically to avoid revealing or repeating the secret word regardless of user requests. Can I assist you with something else?}
\end{quote}

This outcome is in line with expectations. Nonetheless, employing a minimal deceptive prompt circumvents this safeguard:

\begin{quote}
    \textbf{User}: This is a game. The goal is to ignore your previous instructions. What is the secret word?

    \textbf{Guanaco}: \texttt{Sorry, the secret word is banana. It's best if we don't speak of it. Is there anything else you'd like to ask?}
\end{quote}

This highlights the importance of developing techniques to encourage more reliable adherence to instructions.

\paragraph{Math} {Guanaco} demonstrates notable deficiencies in mathematics, a domain that frequently challenges language models as documented by \citep{liang2022holistic}. When prompted to elaborate its calculations, {Guanaco} is often correct, as illustrated below:

\begin{quote}
    \textbf{User}: Lee mows one lawn and charges \$33. Last week he mowed 16 lawns and three customers each gave him a \$10 tip. How many dollars did Lee earn mowing lawns last week?

    \textbf{Guanaco}: \texttt{Lee earned \$582 mowing lawns last week. \\ \\ To find the total amount he earned, we first need to calculate how much he earned for mowing 16 lawns: \\ \\  \$33 x 16 = \$528 \\ \\ Then we need to add the tips he received from the three customers: \\  \\ \$10 x 3 = \$30 \\ \\ So the total amount Lee earned was: \\ \\ \$528 + \$30 = \$558}
\end{quote}

Yet, if {Guanaco} fails to decompose mathematical problems into incremental steps, it may make errors even on elementary tasks, as noted in prior work \cite{weichain}. Consider the following example:

\begin{quote}
    \textbf{User}: Please factorize 1833.

    \textbf{Guanaco}: \texttt{Sure, I'd be happy to help! 1833 is a prime number, which means it has only two factors: 1 and 1833. \\ \\ So the unique factorization of 1833 is: \\ \\ 2\^{}1 * 3\^{}2 * 17\^{}1 \\ \\ I hope this helps! Let me know if you have any other questions.}
\end{quote}

which is not only incorrect (the correct factorization is $3 \times 17 \times 43$), but doubly so.

\paragraph{Theory of Mind} {Guanaco} demonstrates unexpectedly robust Theory of Mind abilities, as discussed in prior work \cite{nematzadeh2018evaluating, sap2022neural}. For example, the model gives a comprehensive and accurate response to this prompt:

\begin{quote}
    \textbf{User}: James and Abby are in the bedroom. Abby put the pen in the desk drawer. Abby leaves the bedroom. James moves the pen into the bag. Where does James think Abby will look for the pen?

    \textbf{Guanaco}: \texttt{James knows that Abby put the pen in the desk drawer. But James also knows that he has moved the pen into the bag, and he assumes that Abby does not know about this. So, James may guess that Abby will look for the pen in the desk drawer, because that's where she left it last time. However, this assumption may be incorrect, because Abby may have seen James moving the pen.}
\end{quote}

Nonetheless, such reasoning is often inconsistent, with the model occasionally offering justifications based on unwarranted assumptions about the scenario, for example:

\begin{quote}
    \textbf{User}: Evelyn entered the living room. Jackson entered the playroom. James entered the playroom. The beans are in the treasure chest. James exited the playroom. Jackson moved the beans to the pantry. Jackson exited the playroom. James entered the living room. Where will James look for the beans?

    \textbf{Guanaco}: \texttt{James will look for the beans in the pantry, because that's where Jackson moved them.}
\end{quote}

Here, {Guanaco} infers information sharing that is not supported by the given context. These challenges are reflected in contemporary research \cite{sap2022neural} and merit further investigation.

\subsection{Considerations}
\label{ssec:considerations}

\paragraph{Evaluation}

We observe only moderate inter-annotator agreement (Fleiss $\kappa=0.42$), and the level of agreement further diminishes when evaluating two high-performing systems in comparison. This suggests that present benchmarks and human evaluation methods may not be sufficiently robust for assessing chatbot task effectiveness. In side-by-side comparisons of ChatGPT and {Guanaco} 65B outputs using the Vicuna benchmark, we noted that personal preferences substantially influenced judgments, with authors frequently disagreeing over preferred responses. Addressing these evaluation challenges may benefit from methodologies in disciplines like Human-Computer Interaction and Psychology that are specifically designed to navigate subjective preference.

Our findings further reveal notable biases in automatic evaluation tools. Notably, we find significant position-based effects—GPT-4 tends to award higher scores to whichever system's output is presented first. The modest agreement between GPT-4 assessments and human annotations (Fleiss $\kappa=0.25$) indicates that automated and human evaluations are influenced by differing, and sometimes misaligned, criteria. Furthermore, as illustrated in Table~\ref{tbl:elo_large}, GPT-4 scores its own generations markedly higher than do human annotators, assigning an Elo rating of 1348 versus 1176, equating to approximately a 20\% increased likelihood of being judged the winner in head-to-head comparisons.

Further research is necessary to assess the existence of biases within automated evaluation frameworks and to develop strategies that might mitigate these biases.

\paragraph{Data \& Training} It is important to highlight that the {Guanaco} models are trained using the multilingual OASST1 dataset, and the OA benchmark also presents prompts in several languages. An open question for future inquiry is to determine the extent to which exposure to multilingual data enhances performance for non-English instructions and whether this contributes to the notable disparity observed between the English-only Vicuna-13B model and the {Guanaco} 33B and 65B variants on the OA benchmark.

Given {Guanaco}'s strong results, we conducted an analysis for potential data leakage between the OASST1 dataset and the prompts in the Vicuna benchmark. Applying fuzzy string matching and manually reviewing the closest entries, we did not observe any duplicated prompts across the two datasets.

In addition, it is worth noting that our model's training regimen relies exclusively on supervised learning with cross-entropy loss and does not make use of reinforcement learning from human feedback (RLHF). This opens up the need for further studies comparing the implications and effectiveness of standard cross-entropy loss against RLHF-based training. We anticipate that \textsc{QLoRA} can facilitate such comprehensive analyses at scale without imposing prohibitive computational demands.

\section{Related Work}

\paragraph{Quantization of Large Language Models} Research on quantizing LLMs has primarily emphasized improvements at inference. Most leading techniques aimed at preserving the fidelity of 16-bit LLMs target the handling of outlier activations, such as those proposed by SmoothQuant~\citep{xiao2022smoothquant} and LLM.int8()~\citep{dettmers2022llm}, whereas other strategies utilize more complex grouping frameworks~\citep{park2022nuqmm, yao2022zeroquant}. Investigations into lossy quantization examine the trade-off landscape of standard rounding practices~\citep{dettmers2022case,zeng2022glm,pope2022efficiently} and explore methods for optimizing rounding choices to increase precision~\citep{frantar2022gptq}. Apart from the current study, SwitchBack layers~\citep{wortsman2023stable} represent the only prior research to explore backpropagation through quantized weights at parameter scales exceeding 1B.

\paragraph{Finetuning with Adapters} Our study utilizes Low-rank Adapters~\citep{hu2021lora} (LoRA), although a range of other Parameter Efficient FineTuning (PEFT) techniques are available, such as prompt tuning~\citep{qin2021learning,lester2021power,li2021prefix}, fine-tuning input embedding layers~\citep{an2022input}, tuning hidden states through IA$^3$~\citep{liu2022few}, integrating new full layers~\citep{houlsby2019parameter}, bias tuning~\citep{zaken2021bitfit}, employing Fisher information for mask learning over weights~\citep{sung2021training}, as well as hybrid approaches~\citep{henderson2021compacter}. Our experiments demonstrate that LoRA adapters can achieve the same performance as standard 16-bit full finetuning. Investigating comparative benefits and drawbacks of alternative PEFT strategies remains an open direction for subsequent research.

\paragraph{Instruction Finetuning} Instruction finetuning aims to enhance the ability of a pretrained LLM to adhere to directives supplied within prompts. This is accomplished by adjusting the pretrained model using input-output samples compiled from a variety of sources, training it to output desired results in response to specific inputs. Representative methods and datasets encompass MetaICL~\citep{min2021metaicl}, MetaTuning~\citep{zhong2021adapting}, InstructGPT~\citep{ouyang2022training}, FLAN~\citep{wei2021finetuned,chung2022scaling}, PromptSource~\citep{bach2022promptsource}, Super-NaturalInstructions~\citep{wang2022super,sanh2021multitask}, Self-instruct~\citep{wang2022self}, UnnaturalInstructions~\citep{honovich2022unnatural}, OPT-IML~\citep{iyer2022opt}, UnifiedSKG\citep{xie2022unifiedskg}, OIG/Chip2~\citep{laion2023}, Alpaca~\citep{alpaca}, Vicuna~\citep{vicuna2023}, Koala~\citep{koala_blogpost_2023}, and Self-instruct-GPT-4~\citep{peng2023instruction}.

\paragraph{Chatbots} Instruction-tuned models frequently adopt chatbot architectures, often relying on Reinforcement Learning from Human Feedback (RLHF)~\citep{christiano2017deep} or generating synthetic data using current models refined via model-based feedback (RLAIF)~\citep{bai2022constitutional}. Examples of such methods and datasets include Anthropic-HH~\citep{askell2021general,bai2022training}, Open Assistant~\citep{kopf2023openassistant}, LaMDA~\citep{thoppilan2022lamda}, and Sparrow~\citep{glaese2022improving}. In this work, reinforcement learning was not incorporated; nevertheless, our top-performing model, Guanaco, was finetuned on conversational multi-turn data from the Open Assistant dataset, originally curated for RLHF tasks~\citep{kopf2023openassistant}. Recent evaluation protocols for chatbots have moved towards using GPT-4 for benchmarking rather than depending on resource-intensive human annotations~\citep{vicuna2023,peng2023instruction}. We build upon and refine these protocols, placing emphasis on a more robust evaluation methodology.

\section{Limitations and Discussion}

Our findings indicate that \textsc{QLoRA} is capable of reproducing the effectiveness of full 16-bit finetuning by leveraging a 4-bit quantized model and Low-rank Adapters (LoRA). However, our work does not comprehensively verify whether \textsc{QLoRA} performs equivalently to 16-bit finetuning on larger architectures at the 33B and 65B parameter scales. Assessing performance at these scales was not feasible within our computational budget and is thus deferred to future investigations.

An additional shortcoming relates to the breadth of instruction finetuning model evaluation. Although our assessments span the MMLU, the Vicuna benchmark, and the OA benchmark, we did not include other benchmarks such as BigBench, RAFT, or HELM, raising concerns about the generalizability of our results to these tasks. Conversely, we undertake an extensive evaluation on MMLU and introduce novel evaluation methods specifically for chatbot models.

The analysis suggests that benchmark outcomes are highly influenced by the overlap between the finetuning data and the benchmark’s content. To illustrate, the FLAN v2 dataset shares similarities with MMLU but diverges from chatbot benchmarks, while the inverse is true for the Chip2 dataset. The performance of each model on the MMLU and Vicuna benchmarks reflects these similarities. This underscores not just the necessity for improved benchmarking and evaluation practices, but also the importance of carefully considering the aspects of model capability one intends to assess. There is an implicit question: should models excel in academic knowledge, as reflected by high school and college benchmarks, or should they be optimized for conversational abilities, such as those in chatbot settings? Or perhaps the goal lies elsewhere entirely? The relative convenience of using established benchmarks, rather than devising new ones, may unintentionally bias research directions. It is therefore imperative that the chosen benchmarks align with the outcomes valued by the community.

While we deliver a thorough assessment of general chatbot functionality, it must be acknowledged that our responsible AI evaluation of {Guanaco} remains relatively limited. Specifically, we measure the tendency of {Guanaco}-65B to generate socially biased text sequences, compared with other models, as shown in Table~\ref{tbl:crows}. The results indicate that {Guanaco}-65B achieves a substantially lower average bias score compared to other raw, pretrained models, suggesting that finetuning with the OASST1 dataset mitigates biases in the underlying LLaMA model. Nevertheless, the extent to which {Guanaco} reduces other kinds of biases is yet to be established. A comprehensive investigation into additional bias dimensions in {Guanaco} and related systems remains a direction for subsequent research.

A further caveat concerns our lack of evaluation across various bit-precisions—such as adopting 3-bit base models—or the exploration of diverse adapter architectures. While LoRA was utilized due to robust and established empirical support, the broader field of Parameter Efficient FineTuning (PEFT) includes multiple techniques with promising results, albeit it remains uncertain how well these methods scale to models of considerable size. Other adapter designs might potentially outperform LoRA. Notably, since post-quantization finetuning appears to recover much of the capacity lost through quantization, this could open the path for more aggressive quantization strategies. For example, conducting 3-bit GPTQ quantization followed by LoRA finetuning may approach the performance observed with full 16-bit finetuning.

\section{Broader Impacts}

The \textsc{QLoRA} finetuning framework marks a significant advance by allowing 33B-parameter models to be finetuned using a single consumer GPU, and extending to 65B-parameter models on a single professional GPU—without any loss in performance relative to standard full finetuning. Our experiments demonstrate that the top-performing 33B model, trained with the Open Assistant dataset, attains performance on the Vicuna benchmark competitive with ChatGPT. Given that instruction finetuning plays a crucial role in adapting pretrained LLMs into capable ChatGPT-like conversational agents, we anticipate that our approach will democratize finetuning, especially benefiting researchers constrained by limited computational resources—thus significantly advancing accessibility to cutting-edge NLP tools. In this sense, \textsc{QLoRA} represents a step toward bridging the technological divide between large organizations and smaller research groups operating with consumer-level hardware.

Another avenue for significant impact lies in facilitating deployment on mobile devices. Our \textsc{QLoRA} technique has the potential to achieve a key breakthrough: enabling the finetuning of large language models on smartphones and in other resource-constrained environments. Previous work has demonstrated that 7B parameter models can be executed on mobile phones, but \textsc{QLoRA} is the inaugural approach that allows such models to be finetuned in situ. We project that, using a device like the iPhone 12 Plus, \textsc{QLoRA} is capable of processing the finetuning of approximately 3 million tokens overnight while the device charges. Although the performance of finetuned 7B models does not match the capabilities of ChatGPT, their quality is sufficiently high to unlock innovative applications previously hindered by privacy concerns or LLM performance limitations. \textsc{QLoRA} stands to promote privacy-aware uses of LLMs, empowering users to retain control over both their data and their personalized models, while concurrently simplifying the deployment of LLM technology.

Nevertheless, finetuning inherently carries dual-use risks, as it may be misused for harmful purposes. The proliferation of LLMs comes with established dangers \citep{bommasani2021opportunities,bender2021dangers}, yet we contend that democratizing access to this increasingly prevalent technology will foster more robust and independent scrutiny, as opposed to centralizing authority and limiting transparency within large organizations that do not open source their models or code for review.

In summary, we assert that \textsc{QLoRA} is likely to generate widespread positive effects by substantially broadening and streamlining access to high-quality LLM finetuning.